\numoptstrue
\yearpaper{2023}{7 May 2023 \textperiodcentered{} Test Booklet Code G6 \textperiodcentered{} Physics Q 1--50 (Section A: 1--35, Section B: 36--50, any 10 of 15) \textperiodcentered{} 180 Marks}{}
\begin{multicols}{2}

\QQ{The work functions of Caesium (Cs), Potassium (K) and Sodium (Na) are $2.14\unit{eV}$, $2.30\unit{eV}$ and $2.75\unit{eV}$ respectively. If incident electromagnetic radiation has an incident energy of $2.20\unit{eV}$, which of these photosensitive surfaces may emit photoelectrons?\ntag{D}{281}{leph203}{8}{ \textperiodcentered{} Kmax=h nu - phi0}{Class XII Part 2 \textperiodcentered{} Ch 11}}
{\opt{Na only}{Cs only}{Both Na and K}{K only}}

\QQ{The net magnetic flux through any closed surface is:\ntag{D}{142}{leph105}{7}{ \textperiodcentered{} net flux through a closed surface is zero}{Class XII Part 1 \textperiodcentered{} Ch 5}}
{\opt{Negative}{Zero}{Positive}{Infinity}}

\QQ{If the galvanometer $G$ does not show any deflection in the circuit shown, the value of $R$ is given by:
\fig{\begin{circuitikz}[scale=0.62,transform shape,font=\Large]
\draw (0,0) to[battery1,l=10\,V,invert] (0,2.4) to[R,l=400\,$\Omega$] (2.4,2.4) to[rmeter,t=G] (4.6,2.4) -- (4.6,1.9);
\draw (4.6,1.9) to[battery1,l_=2\,V] (4.6,0.5) -- (4.6,0) -- (2.4,0) to[short,i<=] (0,0);
\draw (2.4,2.4) to[R,l_=$R$] (2.4,0);
\end{circuitikz}}\ntag{A}{98}{leph103}{18}{}{Class XII Part 1 \textperiodcentered{} Ch 3}}
{\opt{400\,$\Omega$}{200\,$\Omega$}{50\,$\Omega$}{100\,$\Omega$}}

\QQ{A $12\unit{V}$, $60\unit{W}$ lamp is connected to the secondary of a step down transformer, whose primary is connected to ac mains of $220\unit{V}$. Assuming the transformer to be ideal, what is the current in the primary winding?\ntag{D}{195}{leph107}{19}{ \textperiodcentered{} ideal transformer}{Class XII Part 1 \textperiodcentered{} Ch 7}}
{\opt{0.37\,A}{0.27\,A}{2.7\,A}{3.7\,A}}

\QQ{A full wave rectifier circuit consists of two p-n junction diodes, a centre-tapped transformer, capacitor and a load resistance. Which of these components remove the ac ripple from the rectified output?\ntag{D}{339}{leph206}{17}{ \textperiodcentered{} capacitor filter}{Class XII Part 2 \textperiodcentered{} Ch 14}}
{\opt{Load resistance}{A centre-tapped transformer}{p-n junction diodes}{Capacitor}}

\QQ{In a plane electromagnetic wave travelling in free space, the electric field component oscillates sinusoidally at a frequency of $2.0\times10^{10}\unit{Hz}$ and amplitude $48\unit{V\,m^{-1}}$. Then the amplitude of oscillating magnetic field is: (Speed of light in free space $= 3\times10^8\unit{m\,s^{-1}}$)\ntag{D}{207}{leph108}{7}{ \textperiodcentered{} B0=E0/c}{Class XII Part 1 \textperiodcentered{} Ch 8}}
{\opt{$1.6\times10^{-6}\unit{T}$}{$1.6\times10^{-9}\unit{T}$}{$1.6\times10^{-8}\unit{T}$}{$1.6\times10^{-7}\unit{T}$}}

\QQ{A metal wire has mass $(0.4\pm0.002)\unit{g}$, radius $(0.3\pm0.001)\unit{mm}$ and length $(5\pm0.02)\unit{cm}$. The maximum possible percentage error in the measurement of density will nearly be:\ntag{A}{6}{keph101}{6}{}{Class XI Part 1 \textperiodcentered{} Ch 1}}
{\opt{1.4\%}{1.2\%}{1.3\%}{1.6\%}}

\QQ{Light travels a distance $x$ in time $t_1$ in air and $10x$ in time $t_2$ in another denser medium. What is the critical angle for this medium?\ntag{A}{229}{leph201}{9}{}{Class XII Part 2 \textperiodcentered{} Ch 9}}
{\opt{$\sin^{-1}\!\left(\dfrac{10t_1}{t_2}\right)$}{$\sin^{-1}\!\left(\dfrac{t_2}{t_1}\right)$}{$\sin^{-1}\!\left(\dfrac{10t_2}{t_1}\right)$}{$\sin^{-1}\!\left(\dfrac{t_1}{10t_2}\right)$}}

\QQ{An electric dipole is placed at an angle of $30^\circ$ with an electric field of intensity $2\times10^5\unit{N\,C^{-1}}$. It experiences a torque equal to $4\unit{N\,m}$. Calculate the magnitude of charge on the dipole, if the dipole length is $2\unit{cm}$.\ntag{D}{27}{leph101}{27}{ \textperiodcentered{} tau=pE sin theta}{Class XII Part 1 \textperiodcentered{} Ch 1}}
{\opt{2\,mC}{8\,mC}{6\,mC}{4\,mC}}

\QQ{Let a wire be suspended from the ceiling (rigid support) and stretched by a weight $W$ attached at its free end. The longitudinal stress at any point of cross-sectional area $A$ of the wire is:\ntag{D}{168}{keph201}{2}{ \textperiodcentered{} stress=F/A}{Class XI Part 2 \textperiodcentered{} Ch 8}}
{\opt{Zero}{$2W/A$}{$W/A$}{$W/2A$}}

\QQ{In hydrogen spectrum, the shortest wavelength in the Balmer series is $\lambda$. The shortest wavelength in the Brackett series is:\ntag{A}{300}{leph204}{11}{ \textperiodcentered{} series limits from En}{Class XII Part 2 \textperiodcentered{} Ch 12}}
{\opt{$16\lambda$}{$2\lambda$}{$4\lambda$}{$9\lambda$}}

\QQ{The temperature of a gas is $-50^\circ$C. To what temperature the gas should be heated so that the rms speed is increased by 3 times?\ntag{A}{250}{keph205}{7}{}{Class XI Part 2 \textperiodcentered{} Ch 12}}
{\opt{223\,K}{$669^\circ$C}{$3295^\circ$C}{3097\,K}}

\QQ{A football player is moving southward and suddenly turns eastward with the same speed to avoid an opponent. The force that acts on the player while turning is:\ntag{A}{54}{keph104}{6}{}{Class XI Part 1 \textperiodcentered{} Ch 4}}
{\opt{along south-west}{along eastward}{along northward}{along north-east}}

\QQ{The ratio of frequencies of fundamental harmonic produced by an open pipe to that of closed pipe having the same length is:\ntag{D}{292}{keph207}{15}{ \textperiodcentered{} open v/2L, closed v/4L}{Class XI Part 2 \textperiodcentered{} Ch 14}}
{\opt{3:1}{1:2}{2:1}{1:3}}

\QQ{The angular acceleration of a body, moving along the circumference of a circle, is:\ntag{D}{105}{keph106}{14}{ \textperiodcentered{} alpha along the axis}{Class XI Part 1 \textperiodcentered{} Ch 6}}
{\opt{along the axis of rotation}{along the radius, away from centre}{along the radius towards the centre}{along the tangent to its position}}

\QQ{Given below are two statements:\par
\textbf{Statement I}: Photovoltaic devices can convert optical radiation into electricity.\par
\textbf{Statement II}: Zener diode is designed to operate under reverse bias in breakdown region.\par
In the light of the above statements, choose the \textit{most appropriate} answer from the options given below:\ntag{X}{}{}{}{}{Zener diode}}
{\opt{Statement I is incorrect but Statement II is correct.}{Both Statement I and Statement II are correct.}{Both Statement I and Statement II are incorrect.}{Statement I is correct but Statement II is incorrect.}}

\QQ{If $\displaystyle\oint_S\vec E\cdot d\vec S = 0$ over a surface, then:\ntag{D}{31}{leph101}{31}{ \textperiodcentered{} zero net flux}{Class XII Part 1 \textperiodcentered{} Ch 1}}
{\opt{the electric field inside the surface is necessarily uniform.}{the number of flux lines entering the surface must be equal to the number of flux lines leaving it.}{the magnitude of electric field on the surface is constant.}{all the charges must necessarily be inside the surface.}}

\QQ{Resistance of a carbon resistor determined from colour codes is $(22000\pm5\%)\unit{\Omega}$. The colour of third band must be:\ntag{X}{}{}{}{}{resistor colour code}}
{\opt{Yellow}{Red}{Green}{Orange}}

\QQ{The magnetic energy stored in an inductor of inductance $4\unit{\mu H}$ carrying a current of $2\unit{A}$ is:\ntag{D}{169}{leph106}{16}{ \textperiodcentered{} U=LI\textasciicircum{}2/2}{Class XII Part 1 \textperiodcentered{} Ch 6}}
{\opt{$8\unit{\mu J}$}{$4\unit{\mu J}$}{4\,mJ}{8\,mJ}}

\QQ{In a series LCR circuit, the inductance $L$ is $10\unit{mH}$, capacitance $C$ is $1\unit{\mu F}$ and resistance $R$ is $100\unit{\Omega}$. The frequency at which resonance occurs is:\ntag{D}{188}{leph107}{12}{ \textperiodcentered{} omega0=1/sqrt(LC)}{Class XII Part 1 \textperiodcentered{} Ch 7}}
{\opt{1.59\,kHz}{15.9\,rad/s}{15.9\,kHz}{1.59\,rad/s}}

\QQ{The magnitude and direction of the current in the following circuit is
\fig{\begin{circuitikz}[scale=0.62,transform shape,font=\Large]
\draw (0,2) node[left]{A} to[R,l=2\,$\Omega$] (2.2,2) -- (2.7,2); \draw (2.7,2.35) -- (2.7,1.65); \draw (2.95,2.6) -- (2.95,1.4); \node at (2.85,3.0) {10\,V};
\draw (2.95,2) -- (3.9,2) node[circ]{} node[below]{E} -- (4.85,2); \draw (4.85,2.6) -- (4.85,1.4); \draw (5.1,2.35) -- (5.1,1.65); \node at (5.0,3.0) {5\,V};
\draw (5.1,2) to[R,l=1\,$\Omega$] (7.4,2) node[right]{B} -- (7.4,0) node[right]{C} to[R,l=7\,$\Omega$] (0,0) node[left]{D} -- (0,2);
\end{circuitikz}}\ntag{A}{95}{leph103}{15}{}{Class XII Part 1 \textperiodcentered{} Ch 3}}
{\opt{1.5\,A from B to A through E}{0.2\,A from B to A through E}{0.5\,A from A to B through E}{$\dfrac59$\,A from A to B through E}}

\QQ{The minimum wavelength of X-rays produced by an electron accelerated through a potential difference of $V$ volts is proportional to:\ntag{D}{288}{leph203}{15}{ \textperiodcentered{} Exercise 11.1}{Class XII Part 2 \textperiodcentered{} Ch 11}}
{\opt{$V^2$}{$\sqrt V$}{$\dfrac1V$}{$\dfrac1{\sqrt V}$}}

\QQ{The errors in the measurement which arise due to unpredictable fluctuations in temperature and voltage supply are:\ntag{X}{}{}{}{}{types of errors}}
{\opt{Random errors}{Instrumental errors}{Personal errors}{Least count errors}}

\QQ{For Young's double slit experiment, two statements are given below:\par
\textbf{Statement I}: If screen is moved away from the plane of slits, angular separation of the fringes remains constant.\par
\textbf{Statement II}: If the monochromatic source is replaced by another monochromatic source of larger wavelength, the angular separation of fringes decreases.\par
In the light of the above statements, choose the \textit{correct} answer from the options given below:\ntag{A}{266}{leph202}{12}{}{Class XII Part 2 \textperiodcentered{} Ch 10}}
{\opt{Statement I is false but Statement II is true.}{Both Statement I and Statement II are true.}{Both Statement I and Statement II are false.}{Statement I is true but Statement II is false.}}

\QQ{A bullet is fired from a gun at the speed of $280\unit{m\,s^{-1}}$ in the direction $30^\circ$ above the horizontal. The maximum height attained by the bullet is ($g = 9.8\unit{m\,s^{-2}}$, $\sin30^\circ = 0.5$):\ntag{D}{39}{keph103}{13}{ \textperiodcentered{} hm=(v sin theta)\textasciicircum{}2/2g}{Class XI Part 1 \textperiodcentered{} Ch 3}}
{\opt{3000\,m}{2800\,m}{2000\,m}{1000\,m}}

\QQ{A Carnot engine has an efficiency of 50\% when its source is at a temperature $327^\circ$C. The temperature of the sink is:\ntag{D}{238}{keph204}{13}{ \textperiodcentered{} eta=1-T2/T1}{Class XI Part 2 \textperiodcentered{} Ch 11}}
{\opt{$200^\circ$C}{$27^\circ$C}{$15^\circ$C}{$100^\circ$C}}

\QQ{The amount of energy required to form a soap bubble of radius $2\unit{cm}$ from a soap solution is nearly: (surface tension of soap solution $= 0.03\unit{N\,m^{-1}}$)\ntag{A}{194}{keph202}{15}{}{Class XI Part 2 \textperiodcentered{} Ch 9}}
{\opt{$50.1\times10^{-4}\unit{J}$}{$30.16\times10^{-4}\unit{J}$}{$5.06\times10^{-4}\unit{J}$}{$3.01\times10^{-4}\unit{J}$}}

\QQ{The half life of a radioactive substance is 20 minutes. In how much time, the activity of substance drops to $\left(\frac1{16}\right)^{\text{th}}$ of its initial value?\ntag{X}{}{}{}{}{radioactive decay law (half-life)}}
{\opt{80 minutes}{20 minutes}{40 minutes}{60 minutes}}

\QQ{The potential energy of a long spring when stretched by $2\unit{cm}$ is U. If the spring is stretched by $8\unit{cm}$, potential energy stored in it will be:\ntag{D}{81}{keph105}{11}{ \textperiodcentered{} U=kx\textasciicircum{}2/2}{Class XI Part 1 \textperiodcentered{} Ch 5}}
{\opt{16U}{2U}{4U}{8U}}

\QQ{The equivalent capacitance of the system shown in the following circuit is:
\fig{\begin{circuitikz}[scale=0.62,transform shape,font=\Large]
\draw (0,1) node[circ]{} node[left]{A} to[C,l=3\,$\mu$F] (2.4,1) -- (3,1) -- (3,2) to[C,l=3\,$\mu$F] (5,2) -- (5,1) -- (6,1) node[circ]{} node[right]{B};
\draw (3,1) -- (3,0) to[C,l_=3\,$\mu$F] (5,0) -- (5,1);
\end{circuitikz}}\ntag{D}{72}{leph102}{28}{ \textperiodcentered{} series and parallel}{Class XII Part 1 \textperiodcentered{} Ch 2}}
{\opt{9\,$\mu$F}{2\,$\mu$F}{3\,$\mu$F}{6\,$\mu$F}}

\QQ{A vehicle travels half the distance with speed $\vartheta$ and the remaining distance with speed $2\vartheta$. Its average speed is:\ntag{A}{21}{keph102}{9}{}{Class XI Part 1 \textperiodcentered{} Ch 2}}
{\opt{$\dfrac{3\vartheta}4$}{$\dfrac\vartheta3$}{$\dfrac{2\vartheta}3$}{$\dfrac{4\vartheta}3$}}

\QQ{The ratio of radius of gyration of a solid sphere of mass $M$ and radius $R$ about its own axis to the radius of gyration of the thin hollow sphere of same mass and radius about its axis is:\ntag{A}{115}{keph106}{24}{}{Class XI Part 1 \textperiodcentered{} Ch 6}}
{\opt{5:2}{3:5}{5:3}{2:5}}

\QQ{Two bodies of mass $m$ and $9m$ are placed at a distance $R$. The gravitational potential on the line joining the bodies where the gravitational field equals zero, will be ($G$ = gravitational constant):\ntag{A}{136}{keph107}{10}{ \textperiodcentered{} Example 7.4 setting}{Class XI Part 1 \textperiodcentered{} Ch 7}}
{\opt{$-\dfrac{20Gm}R$}{$-\dfrac{8Gm}R$}{$-\dfrac{12Gm}R$}{$-\dfrac{16Gm}R$}}

\QQ{The venturi-meter works on:\ntag{A}{188}{keph202}{9}{ \textperiodcentered{} venturimeter (Bernoulli)}{Class XI Part 2 \textperiodcentered{} Ch 9}}
{\opt{The principle of perpendicular axes}{Huygen's principle}{Bernoulli's principle}{The principle of parallel axes}}

\QQ{An ac source is connected to a capacitor C. Due to decrease in its operating frequency:\ntag{A}{185}{leph107}{9}{}{Class XII Part 1 \textperiodcentered{} Ch 7}}
{\opt{capacitive reactance remains constant}{capacitive reactance decreases.}{displacement current increases.}{displacement current decreases.}}

\QQ{The radius of inner most orbit of hydrogen atom is $5.3\times10^{-11}\unit{m}$. What is the radius of third allowed orbit of hydrogen atom?\ntag{D}{305}{leph204}{16}{ \textperiodcentered{} Exercise 12.7}{Class XII Part 2 \textperiodcentered{} Ch 12}}
{\opt{4.77\,\AA}{0.53\,\AA}{1.06\,\AA}{1.59\,\AA}}

\QQ{The resistance of platinum wire at $0^\circ$C is $2\unit{\Omega}$ and $6.8\unit{\Omega}$ at $80^\circ$C. The temperature coefficient of resistance of the wire is:\ntag{D}{90}{leph103}{10}{ \textperiodcentered{} rho=rho0[1+alpha(T-T0)]}{Class XII Part 1 \textperiodcentered{} Ch 3}}
{\opt{$3\times10^{-1}\,{}^\circ\mathrm C^{-1}$}{$3\times10^{-4}\,{}^\circ\mathrm C^{-1}$}{$3\times10^{-3}\,{}^\circ\mathrm C^{-1}$}{$3\times10^{-2}\,{}^\circ\mathrm C^{-1}$}}

\QQ{The net impedance of circuit (as shown in figure) will be:
\fig{\begin{circuitikz}[scale=0.62,transform shape,font=\Large]
\draw (0,2) to[L,l_={$\frac{50}{\pi}$\,mH}] (2,2) to[C,l_={$\frac{10^3}{\pi}\,\mu$F}] (3.6,2) to[R,l=10\,$\Omega$] (5.4,2) -- (5.4,0) to[sV] (0,0) -- (0,2); \node at (2.7,-0.8) {220\,V, 50\,Hz};
\end{circuitikz}}\ntag{A}{187}{leph107}{11}{}{Class XII Part 1 \textperiodcentered{} Ch 7}}
{\opt{25\,$\Omega$}{$10\sqrt2$\,$\Omega$}{15\,$\Omega$}{$5\sqrt5$\,$\Omega$}}

\QQ{For the following logic circuit, the truth table is:
\fig{\begin{circuitikz}[scale=0.7,transform shape]
\draw (0,1) node[not port](N1){}; \draw (0,-0.6) node[not port](N2){}; \draw (2.2,0.2) node[nand port](G){};
\draw (N1.in) -- ++(-0.5,0) node[left]{A}; \draw (N2.in) -- ++(-0.5,0) node[left]{B};
\draw (N1.out) -| (G.in 1); \draw (N2.out) -| (G.in 2); \draw (G.out) -- ++(0.4,0) node[right]{Y};
\end{circuitikz}}\ntag{X}{}{}{}{}{digital electronics (logic gates)}}
{\opt{\ttab{0&0&0}{0&1&0}{1&0&0}{1&1&1}}{\ttab{0&0&1}{0&1&1}{1&0&1}{1&1&0}}{\ttab{0&0&0}{0&1&1}{1&0&1}{1&1&1}}{\ttab{0&0&1}{0&1&0}{1&0&1}{1&1&0}}}

\QQ{10 resistors, each of resistance $R$ are connected in series to a battery of emf $E$ and negligible internal resistance. Then those are connected in parallel to the same battery, the current is increased $n$ times. The value of $n$ is:\ntag{A}{98}{leph103}{18}{}{Class XII Part 1 \textperiodcentered{} Ch 3}}
{\opt{1000}{10}{100}{1}}

\QQ{Calculate the maximum acceleration of a moving car so that a body lying on the floor of the car remains stationary. The coefficient of static friction between the body and the floor is 0.15 ($g = 10\unit{m\,s^{-2}}$).\ntag{D}{61}{keph104}{13}{ \textperiodcentered{} Example 4.7}{Class XI Part 1 \textperiodcentered{} Ch 4}}
{\opt{$50\unit{m\,s^{-2}}$}{$1.2\unit{m\,s^{-2}}$}{$150\unit{m\,s^{-2}}$}{$1.5\unit{m\,s^{-2}}$}}

\QQ{The $x$-$t$ graph of a particle performing simple harmonic motion is shown in the figure. The acceleration of the particle at $t = 2\unit{s}$ is:
\fig{\begin{tikzpicture}[xscale=0.4,yscale=0.8]
\draw[->] (0,0) -- (9.5,0) node[right,lbl]{$t$ (s)}; \draw[->] (0,-1.3) -- (0,1.4) node[above,lbl]{$x$ (m)};
\draw[domain=0:8,samples=60] plot (\x,{sin(45*\x)});
\draw[densely dashed] (0,1) node[left,lbl]{1} -- (2,1) -- (2,0) node[below,lbl]{2} (0,-1) node[left,lbl]{$-1$} -- (6,-1) -- (6,0) node[above,lbl]{6};
\node[lbl,above] at (4,0) {4}; \node[lbl,below] at (8,0) {8};
\end{tikzpicture}}\ntag{A}{266}{keph206}{8}{}{Class XI Part 2 \textperiodcentered{} Ch 13}}
{\opt{$-\dfrac{\pi^2}{16}\unit{m\,s^{-2}}$}{$\dfrac{\pi^2}8\unit{m\,s^{-2}}$}{$-\dfrac{\pi^2}8\unit{m\,s^{-2}}$}{$\dfrac{\pi^2}{16}\unit{m\,s^{-2}}$}}

\QQ{A satellite is orbiting just above the surface of the earth with period $T$. If $d$ is the density of the earth and $G$ is the universal constant of gravitation, the quantity $\dfrac{3\pi}{Gd}$ represents:\ntag{A}{137}{keph107}{11}{}{Class XI Part 1 \textperiodcentered{} Ch 7}}
{\opt{$\sqrt T$}{$T$}{$T^2$}{$T^3$}}

\QQ{A very long conducting wire is bent in a semi-circular shape from A to B as shown in figure. The magnetic field at point P for steady current configuration is given by:
\fig{\begin{tikzpicture}[scale=0.8]
\draw[densely dashed] (-2.6,0.9) -- (-1.4,0.9); \draw[->] (-1.4,0.9) -- (-0.6,0.9); \draw (-0.6,0.9) -- (0,0.9) node[right,lbl]{A}; \node[lbl] at (-0.9,1.2) {$i\rightarrow$};
\draw (0,0.9) arc (90:-90:0.9); \draw[->] (0.64,-0.64) -- (0.6,-0.67);
\draw (0,-0.9) node[right,lbl]{B} -- (-0.6,-0.9); \draw[->] (-0.6,-0.9) -- (-1.4,-0.9); \draw[densely dashed] (-1.4,-0.9) -- (-2.6,-0.9); \node[lbl] at (-0.9,-0.6) {$i\leftarrow$};
\fill (0,0) circle (1.3pt) node[right,lbl]{P}; \draw[->] (0,0) -- (130:0.9) node[pos=0.6,above right,lbl]{$R$};
\end{tikzpicture}}\ntag{X}{}{}{}{}{magnetic field of a finite straight wire}}
{\opt{$\dfrac{\mu_0i}{4R}\left[1-\dfrac2\pi\right]$ pointed into the page}{$\dfrac{\mu_0i}{4R}$ pointed into the page}{$\dfrac{\mu_0i}{4R}$ pointed away from the page}{$\dfrac{\mu_0i}{4R}\left[1-\dfrac2\pi\right]$ pointed away from page}}

\QQ{In the figure shown here, what is the equivalent focal length of the combination of lenses (Assume that all layers are thin)?
\fig{\begin{tikzpicture}[scale=0.8]
\fill[pattern=horizontal lines,pattern color=gray] (0,0) rectangle (1.8,2.4); \draw (0,0) rectangle (1.8,2.4);
\fill[white] (0.9,2.4) .. controls (0.35,1.8) and (0.35,0.6) .. (0.9,0) .. controls (1.45,0.6) and (1.45,1.8) .. (0.9,2.4);
\draw (0.9,2.4) .. controls (0.35,1.8) and (0.35,0.6) .. (0.9,0) .. controls (1.45,0.6) and (1.45,1.8) .. (0.9,2.4);
\node[lbl] at (0.7,0.5) {$R_1$}; \node[lbl] at (1.15,0.35) {$R_2$};
\draw[->] (-0.8,2.6) node[left,lbl]{$n_1 = 1.5$} -- (0.9,2.6) -- (0.9,2.2); \draw[->] (-0.8,-0.2) node[left,lbl]{$n_2 = 1.6$} -- (0.3,-0.2) -- (0.3,0.25);
\node[lbl,right] at (1.9,1.2) {$R_1 = R_2 = 20\unit{cm}$};
\end{tikzpicture}}\ntag{A}{238}{leph201}{18}{}{Class XII Part 2 \textperiodcentered{} Ch 9}}
{\opt{$-50\unit{cm}$}{40\,cm}{$-40\unit{cm}$}{$-100\unit{cm}$}}

\QQ{Two thin lenses are of same focal lengths ($f$), but one is convex and the other one is concave. When they are placed in contact with each other, the equivalent focal length of the combination will be:\ntag{D}{237}{leph201}{17}{ \textperiodcentered{} 1/f=1/f1+1/f2}{Class XII Part 2 \textperiodcentered{} Ch 9}}
{\opt{Infinite}{Zero}{$f/4$}{$f/2$}}

\QQ{A wire carrying a current $I$ along the positive $x$-axis has length $L$. It is kept in a magnetic field $\vec B = (2\hat i+3\hat j-4\hat k)\unit{T}$. The magnitude of the magnetic force acting on the wire is:\ntag{D}{110}{leph104}{4}{ \textperiodcentered{} F=I l x B}{Class XII Part 1 \textperiodcentered{} Ch 4}}
{\opt{$\sqrt3\,IL$}{$3IL$}{$\sqrt5\,IL$}{$5IL$}}

\QQ{A bullet from a gun is fired on a rectangular wooden block with velocity $u$. When bullet travels $24\unit{cm}$ through the block along its length horizontally, velocity of bullet becomes $\dfrac u3$. Then it further penetrates into the block in the same direction before coming to rest exactly at the other end of the block. The total length of the block is:\ntag{A}{17}{keph102}{5}{}{Class XI Part 1 \textperiodcentered{} Ch 2}}
{\opt{30\,cm}{27\,cm}{24\,cm}{28\,cm}}

\QQ{An electric dipole is placed as shown in the figure. The electric potential (in $10^2\unit{V}$) at point P due to the dipole is ($\varepsilon_0$ = permittivity of free space and $\dfrac1{4\pi\varepsilon_0} = K$):
\fig{\begin{tikzpicture}[scale=0.55]
\fill (-3,0) circle (2.5pt) node[below left,lbl]{$-q$}; \fill (3,0) circle (2.5pt) node[below right,lbl]{$+q$}; \draw (-3,0) -- (3,0);
\draw (0,-0.2) -- (0,0.2) node[above,lbl]{0}; \fill (5,0) circle (2.5pt) node[above right,lbl]{P};
\draw[<->] (-3,-0.6) -- (0,-0.6) node[midway,below,lbl]{3\,cm}; \draw[<->] (0,-0.6) -- (3,-0.6) node[midway,below,lbl]{3\,cm};
\draw (0,0.8) -- (5,0.8) node[midway,above,lbl]{5\,cm};
\end{tikzpicture}}\ntag{A}{51}{leph102}{7}{}{Class XII Part 1 \textperiodcentered{} Ch 2}}
{\opt{$\left(\dfrac83\right)qK$}{$\left(\dfrac38\right)qK$}{$\left(\dfrac58\right)qK$}{$\left(\dfrac85\right)qK$}}

\QQ{A horizontal bridge is built across a river. A student standing on the bridge throws a small ball vertically upwards with a velocity $4\unit{m\,s^{-1}}$. The ball strikes the water surface after $4\unit{s}$. The height of bridge above water surface is (Take $g = 10\unit{m\,s^{-2}}$):\ntag{D}{18}{keph102}{6}{ \textperiodcentered{} Example 2.3}{Class XI Part 1 \textperiodcentered{} Ch 2}}
{\opt{68\,m}{56\,m}{60\,m}{64\,m}}

%%BEGIN-ANALYSIS
\end{multicols}
\par\needspace{16\baselineskip}\vspace{14pt}
\begin{center}{\pgfornament[width=3.2cm,color=gilt]{87}}\par
{\large\color{navy}\scshape NEET 2023 \textperiodcentered{} Answer Key}\end{center}
{\small\setlength{\tabcolsep}{1pt}\renewcommand{\arraystretch}{1.5}
\begin{tabularx}{\linewidth}{@{}*{10}{>{\centering\arraybackslash}X}@{}}\toprule
\textbf{1}\,(2) & \textbf{2}\,(2) & \textbf{3}\,(4) & \textbf{4}\,(2) & \textbf{5}\,(4) & \textbf{6}\,(4) & \textbf{7}\,(4) & \textbf{8}\,(1) & \textbf{9}\,(1) & \textbf{10}\,(3)\\
\textbf{11}\,(3) & \textbf{12}\,(3) & \textbf{13}\,(4) & \textbf{14}\,(3) & \textbf{15}\,(1) & \textcolor{gray}{\textbf{16}\,(2)$^{\dagger}$} & \textbf{17}\,(2) & \textcolor{gray}{\textbf{18}\,(4)$^{\dagger}$} & \textbf{19}\,(1) & \textbf{20}\,(1)\\
\textbf{21}\,(3) & \textbf{22}\,(3) & \textcolor{gray}{\textbf{23}\,(1)$^{\dagger}$} & \textbf{24}\,(4) & \textbf{25}\,(4) & \textbf{26}\,(2) & \textbf{27}\,(4) & \textcolor{gray}{\textbf{28}\,(1)$^{\dagger}$} & \textbf{29}\,(1) & \textbf{30}\,(2)\\
\textbf{31}\,(4) & \textbf{32}\,--$^{*}$ & \textbf{33}\,(4) & \textbf{34}\,(3) & \textbf{35}\,(4) & \textbf{36}\,(1) & \textbf{37}\,(4) & \textbf{38}\,(4) & \textcolor{gray}{\textbf{39}\,(3)$^{\dagger}$} & \textbf{40}\,(3)\\
\textbf{41}\,(4) & \textbf{42}\,(1) & \textbf{43}\,(3) & \textcolor{gray}{\textbf{44}\,(1)$^{\dagger}$} & \textbf{45}\,(4) & \textbf{46}\,(1) & \textbf{47}\,(4) & \textbf{48}\,(2) & \textbf{49}\,(2) & \textbf{50}\,(4)\\\bottomrule\end{tabularx}}\par\smallskip
{\footnotesize Official NTA final key (Test Booklet Code G6); every answer was also solved and checked by the editor. (1)--(4) are the printed options. {\color{gray}Grey$^{\dagger}$}: deleted portion (answer given for completeness). --: question dropped by NTA (marks given to all). $^{*}$: see note.\par
\par\smallskip\textbf{Q32.} Dropped by NTA. The ratio is $\sqrt{2/5}:\sqrt{2/3}=\sqrt3:\sqrt5$, which is not among the options.
}\par
\clearpage
\begin{center}{\pgfornament[width=5cm,color=gilt]{87}}\par
{\fontsize{20}{24}\selectfont\color{navy}\scshape NEET 2023 \textperiodcentered{} Paper Analysis}\par
{\color{gilt}\pgfornament[width=5cm,color=gilt,symmetry=h]{87}}\end{center}
\vspace{4pt}
\begin{multicols}{2}
\textbf{\color{crimson}1. Direct vs Applied}\par\smallskip
{\footnotesize Checked question by question against the four 2026-27 NCERT books (XI Part 1 \& 2, XII Part 1 \& 2). Out of 50, 44 questions are in the present syllabus; percentages are of these.}\par\smallskip
\begin{tabular}{@{}lrr@{}}\toprule
Direct from NCERT & \textbf{23} & 52\%\\
\quad text, summary, table or figure & 19 &\\
\quad NCERT solved Example & 2 &\\
\quad NCERT Exercise & 2 &\\
Applied (concept in NCERT) & \textbf{21} & 48\%\\\midrule
Analysed & 44 &\\
Deleted portion (not analysed) & 6 &\\\midrule
\textbf{Total} & \textbf{50} &\\\bottomrule\end{tabular}\par\smallskip
{\footnotesize Direct: Q 1, 2, 4, 5, 6, 9, 10, 14, 15, 17, 19, 20, 22, 25, 26, 29, 30, 36, 37, 41, 46, 47, 50.\\ From NCERT Examples: Q 41 (Example 4.7), 50 (Example 2.3).\\ From NCERT Exercises: Q 22 (Exercise 11.1), 36 (Exercise 12.7).}\par\medskip
\textbf{\color{crimson}2. Difficulty level}\par\smallskip
\begin{tabular}{@{}lr@{}}\toprule
Easy & 25\\ Moderate & 18\\ Hard & 1\\\bottomrule\end{tabular}\par\smallskip
{\footnotesize Hard: Q 45. The 2023 paper was easy: just over half the in-syllabus questions are one-step uses of NCERT formulas, examples and exercises. Class XI: 19 questions, Class XII: 25 (analysed questions only). (Difficulty is the editor's assessment.)}\par\medskip
\columnbreak
\textbf{\color{crimson}3. Chapter-wise distribution}\par\smallskip
{\footnotesize\begin{tabular}{@{}p{4.7cm}c@{\,}r@{\,}r@{\,}r@{}}\toprule
Chapter & Cl. & Qs & Dir & App\\\midrule
Units \& Measurement & XI & 1 & 0 & 1\\
Motion in a Straight Line & XI & 3 & 1 & 2\\
Motion in a Plane & XI & 1 & 1 & 0\\
Laws of Motion & XI & 2 & 1 & 1\\
Work, Energy \& Power & XI & 1 & 1 & 0\\
Rotational Motion & XI & 2 & 1 & 1\\
Gravitation & XI & 2 & 0 & 2\\
Mech.\ Prop.\ of Solids & XI & 1 & 1 & 0\\
Mech.\ Prop.\ of Fluids & XI & 2 & 0 & 2\\
Thermodynamics & XI & 1 & 1 & 0\\
Kinetic Theory & XI & 1 & 0 & 1\\
Oscillations & XI & 1 & 0 & 1\\
Waves & XI & 1 & 1 & 0\\
Electric Charges \& Fields & XII & 2 & 2 & 0\\
Electrostatic Potential \& Capacitance & XII & 2 & 1 & 1\\
Current Electricity & XII & 4 & 1 & 3\\
Moving Charges \& Magnetism & XII & 1 & 1 & 0\\
Magnetism \& Matter & XII & 1 & 1 & 0\\
Electromagnetic Induction & XII & 1 & 1 & 0\\
Alternating Current & XII & 4 & 2 & 2\\
Electromagnetic Waves & XII & 1 & 1 & 0\\
Ray Optics & XII & 3 & 1 & 2\\
Wave Optics & XII & 1 & 0 & 1\\
Dual Nature & XII & 2 & 2 & 0\\
Atoms & XII & 2 & 1 & 1\\
Semiconductor Electronics & XII & 1 & 1 & 0\\
\bottomrule\end{tabular}}\par\medskip
\textbf{\color{crimson}4. Teaching takeaways}\par\smallskip
{\footnotesize\begin{itemize}[leftmargin=1em,nosep]
\item NCERT Examples 2.3 and 4.7 and Exercises 11.1 and 12.7 were asked with the same idea or the same numbers.
\item One-line formulas ($\omega_0=1/\sqrt{LC}$, $B_0=E_0/c$, $\tfrac12LI^2$, Carnot efficiency) gave quick marks.
\item Six questions were on deleted portions (Zener diode, colour code, types of errors, half-life, logic gates, finite wire).
\item Q32 (radius of gyration) was dropped by NTA because no option is correct.
\end{itemize}}\par\medskip
\textbf{\color{crimson}5. Deleted portions}\par\smallskip
\textbf{6} of the 50 questions are on topics removed from the 2026-27 NCERT (rationalised syllabus). They are marked \textsc{Deleted Portion} and are not analysed: Zener diode (Q 16); resistor colour code (Q 18); types of errors (Q 23); radioactive decay law (half-life) (Q 28); digital electronics (logic gates) (Q 39); magnetic field of a finite straight wire (Q 44).\par
\end{multicols}
%%END-ANALYSIS


